% ============================================================
% 二、模型假设与符号说明
% ============================================================
\section{模型假设}

为使问题在合理范围内可解，我们做出以下基本假设：

（以段落形式逐条论述每个假设，说明假设的依据和合理性。
不要使用 enumerate 或 itemize 列表。
例如："假设一：我们假设观测数据中不存在系统性测量偏差。
该假设基于...的事实，在实际工况中..."
每个假设一段，通常 5--8 条假设。
随着各子问题的推进，Writer 可在此追加新引入的假设。）

\section{符号说明}

\begin{table}[H]
\centering
\caption{主要符号说明}
\begin{tabular}{cl}
\toprule
\textbf{符号} & \textbf{说明} \\
\midrule
$x$ & 示例变量（Writer 根据实际模型替换） \\
$y$ & 示例变量 \\
\bottomrule
\end{tabular}
\end{table}
